\setsubtitle{
    author : blame\\
    最高分 : TODO
}
\section{E. 分蛋糕4}

\begin{frame}{\btitle{題目敘述}}
    題目太長了

    懶得打，自己去看題目。
\end{frame}

\begin{frame}{\btitle{子任務}}
    \begin{enumerate}
        \item{$2 \le n \le 50, 1 \le Q \le 1000$}
        \item{$2 \le n \le 200$}
        \item{$2 \le n \le 800$}
        \item{無額外限制}
    \end{enumerate}
\end{frame}

\begin{frame}{\btitle{$2 \le n \le 50, 1 \le Q \le 1000$}}
    你可以把整個陣列原封不動傳給 decoder

    然後讓 decoder 自己算+回答
\end{frame}

\begin{frame}{\btitle{$2 \le n \le 200$}}
    區間 DP 做完後，傳給 decoder

    複雜度 $O(n^3)$
\end{frame}

\begin{frame}{\btitle{$2 \le n \le 800$}}
    發現這題的 $cost$ 其實有性質

    可以用 knuth 做到 $O(n^2)$ 算

    然後樸實無華的傳給 decoder 回答
\end{frame}

\begin{frame}{\btitle{無額外限制}}
    這邊講做到 $\lvert S \rvert$ 為 $2 \times C^{n}_{2}$ 的作法

    先用 knuth 算出答案後

    然後發現對於每一列的 $k$ 其實具有遞增性

    所以你只要用 `0` 作為斷點，計算每個斷點中有多少個 `1`

    代表著 $k$ 相對於之前那個增加多少
\end{frame}

\begin{frame}{\btitle{無額外限制}}
    這邊講做到 $\lvert S \rvert$ 為 $2850423$ 的作法

    你發現一列一列傳沒有完整用到 knuth 的性質

    所以你用一個斜線一個斜線的方式傳(傳的過程可以用一個大整數代表著整列的狀態)

    然後數學計算一下發現是好的。\pause

    假如你還要做到更佳，可以去問 brinton
\end{frame}

